
\section{Setup} \label{sec:setup}

We consider an agent with an expected utility function facing a consumption problem
over a sequence of periods. At each period $t$, the agent chooses consumption $c_t$,
and perfectly observes state variable $s_t$. The utility function is a function of
$c_t$ and $s_t$ and therefore the agent perfectly observes their payoff $u(c_t, s_t)$ at
each period. The agent discounts future utility by a factor $\beta \in (0,1)$.

\begin{equation}
    V^*(s_t) = \max_{c_t \in B(s_t)} \left\{ u(c_t, s_t) + \beta \mathbb{E} \left[ V^*(s_{t+1}) \right] \right\}
\end{equation}

This is the Bellman equation for the agent's state value function $V^*(s_t)$, where
$B(s_t)$ is the agent's budget set at state $s_t$. This is directly linked to the
agent's action value function $Q^*(s_t, c_t)$ and optimal policy function as follows.\footnote{Here we are using 
$*$ to indicate the state and action value functions under the optimal policy,
there is no difference in the formation of the equations under different policies. 
Notationally, it would be indicated as $V^\pi$ and $Q^\pi$ for a policy $\pi$.}.

\begin{proposition}[Relationships Between $Q^*$, $V^*$, and $\pi^*$] \label{prop:bellman_relationships}
The optimal state value function $V^*(s)$, the optimal action value function $Q^*(s, c)$, and the optimal policy $\pi^*(s)$ are related as follows:
\begin{align*}
    \pi^*(s) &= \arg\max_{c \in B(s)} Q^*(s, c), \\
    V^*(s) &= \max_{c \in B(s)} Q^*(s, c) = Q^*(s, \pi^*(s)), \\
    Q^*(s, c) &= u(c, s) + \beta \mathbb{E}[V^*(s') \mid s, c], \\
    V^*(s) &= u(\pi^*(s), s) + \beta \mathbb{E}[V^*(s') \mid s, \pi^*(s)].
\end{align*}
\end{proposition}

There is a stochastic Markov transition function $\mathcal{F}(c_t, s_t, \epsilon_{t+1})$
where $\epsilon_{t+1} \sim \mathcal{G}(\epsilon_{t+1})$ is an i.i.d. shock that perturbs
the transition from $s_t$ to $s_{t+1}$.

\subsection{Agent's Beliefs}

The agent believes there to be some unknown optimal policy $\pi^*$ that produces
optimal state action value function $Q^*(s_t, c_t)$. As the agent has not learned this policy
yet, but they do form a prior over the set of possible action value functions $\mathcal{Q}$, which
are modeled as Gaussian Processes (GPs).

\begin{equation}
    Q^*(c,s) \sim GP\left(\hat{Q}_0(c,s), \hat{\Sigma_0}((s,c), (s', c'))\right)
\end{equation}

The mean function is $\hat{Q}_0(c,s)$ and the covariance function is $\hat{\Sigma_0}((s,c), (s', c'))$.
The use of GPs becomes useful in that they allow the agent to have correlated beliefs over
similar state-action pairs, allowing for generalization across the space $\mathcal{Q}$. And,
as new observations are made, the agent can update their beliefs using Bayesian updating, 
which results in a new GP posterior over $\mathcal{Q}$. There are some useful properties 
of GPs that will list out here for reference.

\begin{proposition}[Closure Under Conditioning]
    If $f \sim GP(\mu, k)$ and we have observations $\mathcal{D} = \{(x_i, y_i)\}_{i=1}^n$ where $y_i = f(x_i) + \epsilon_i$ with $\epsilon_i \sim \mathcal{N}(0, \sigma^2)$.
    The posterior distribution $f | \mathcal{D} \sim GP(\mu_n, k_n)$ where
\begin{align*}
\mu_n(x) &= \mu(x) + k(x, X_n)(k(X_n, X_n) + \sigma^2 I)^{-1}(Y_n - \mu(X_n)) \\
        k_n(x, x') &= k(x, x') - k(x, X_n)(k(X_n, X_n) + \sigma^2 I)^{-1}k(X_n, x')
\end{align*}
    where $X_n = (x_1, \ldots, x_n)^\top$ is the matrix of observed inputs and $Y_n = (y_1, \ldots, y_n)^\top$ is the vector of observed outputs.
\end{proposition}

\begin{proposition}[Finite-Dimensional Marginals Are Gaussian]
    For any finite collection of points $\{x_1, \ldots, x_k\}$, the joint distribution $(f(x_1), \ldots, f(x_k))$ is multivariate normal. This implies that predictions at new points have joint Gaussian distributions.
\end{proposition}

\begin{proposition}[Linearity of the Mean]
    If $f \sim GP(\mu, k)$, then for any linear functional $L$:
    \begin{equation*}
        L[f] \sim \mathcal{N}(L[\mu], L[k])
    \end{equation*}
    This extends to integration against measures. For instance, if the agent needs $\mathbb{E}_{s'}[V^*(s')] = \mathbb{E}_{s'}[\max_c Q^*(s', c)]$, the inner expectation over the transition kernel is tractable.
    This is because the expectation operator defines a bounded linear functional on $L^1(\Omega, \mathcal{F}, P)$.
    though the $\max$ operator introduces nonlinearity.
\end{proposition}

\subsection{Learning Process}

Given the agent's prior is conveniently expressed as a gaussian process over action value functions,
the agent performs Bayesian updating as new observations are made. The agent updates her beliefs 
over $\mathcal{Q}$ through two mechanisms: experience  and reasoning.\footnote{
    Given the agent is endowed with a belief over the action value function $Q^*$, they can derive
    implied beliefs over the optimal policy $\pi^*$ and state value function $V^*$ using the relationships
    stated in Proposition \ref{prop:bellman_relationships}. To compute the agent's prior belief over policies 
    $\pi$, we need to calculate the distribution of the $\arg\max$ of a GP. However, this is not analytically 
    tractable: if $f \sim GP(\mu, k)$, then $\arg\max_x f(x)$ is a random element of the input space whose 
    distribution depends on the entire sample path structure and has no closed form.
}
The experience mechanism involves updating beliefs through direct observation of realized outcomes 
$(s_t, c_t, u_t, s_{t+1})$ from actions taken. The reasoning mechanism involves updating 
beliefs through mental simulation---evaluating counterfactual actions using the agent's model of transitions $\mathcal{F}$ and 
payoffs $u$ without taking those actions.

\begin{definition}[The Empirical-Reasoning Learning Cycle] \label{def:learning_cycle}
    At each period $t$, the agent immediately receives signal $u(c_{t-1}, s_{t-1})$ from the previous period's action.
    She then first performs an experience-based update to form intermediate beliefs $\hat{Q}^E_t, \hat{\Sigma}^E_t$.
    Which has a corresponding optimal policy $\pi^E_t$ and state value function $V^E_t$.
    Next, she performs a reasoning-based update using mental simulation under the new policy $\pi^E_t$ to form
    reasoned beliefs $\hat{Q}^R_t, \hat{\Sigma}^R_t$, with corresponding $\pi^R_t$ and $V^R_t$.
    Finally, she chooses action $c_t = \pi^R_t(s_t)$ under the reasoned policy.
\end{definition}


\subsubsection{Experience-Based Learning}

The agent takes a consumption action $c_{t-1}$ in state $s_{t-1}$, receives
payoff $u_{t-1} = u(c_{t-1}, s_{t-1})$, and transitions to state $s_t$, based on
the Markov transition function $\mathcal{F}(c_{t-1}, s_{t-1}, \epsilon_t)$. They can
now look back and form a temporal difference (TD) approximation. By rearranging the
Bellman equation we can write the realized utility as a function of the action value function.
\begin{equation}
\underbrace{u(c_{t-1}, s_{t-1})}_{\text{observed}} = \underbrace{Q^*(c_{t-1}, s_{t-1})}_{\text{unknown, GP prior}} - \beta \underbrace{\mathbb{E}_{t-1}[Q^*(s_t, \pi^*(s_t))]}_{\text{requires integration over } \mathcal{F}}
\end{equation}

Since the agent does not know $Q^*$ or the transition function $\mathcal{F}$ 
\textit{ex ante}, they cannot compute the Bellman expectation analytically. In the 
experience-based learning phase, the agent replaces the theoretical expectation over 
all possible future states with the single realized transition to $s_t$. Rather than 
performing a simulation of potential outcomes, the agent learns directly from the 
empirical evidence of the outcome. By observing the arrived state $s_t$ and identifying 
their current best guess of the optimal action, $\tilde{c}_t = \arg\max_{c} \hat{Q}_{t-1}(c, s_t)$, 
the agent constructs a Temporal Difference (TD) signal. This signal serves as a 
sample-based realization of the Bellman equation, which the agent then uses to 
perform a Bayesian update on the GP prior.

\begin{definition}[Empirical Learning Update]
    Upon arriving at period $t$, the agent observes the realized utility $\eta^E_t = u(c_{t-1}, s_{t-1})$ and the transition to state $s_t$. Given the prior beliefs $(\hat{Q}_{t-1}, \hat{\Sigma}_{t-1})$, the agent treats the experience as a signal $\eta^E_t$ of the true $Q^*$ function. The posterior beliefs $(\hat{Q}^E_t, \hat{\Sigma}^E_t)$ are formed via a Bayesian update:

    \begin{equation}
        \hat{Q}^E_t(c,s) = \hat{Q}_{t-1}(c,s) + \alpha^E_t(c,s) \underbrace{\left[ \eta^E_t + \beta \hat{Q}_{t-1}(\tilde{c}_t, s_t) - \hat{Q}_{t-1}(c_{t-1}, s_{t-1}) \right]}_{\text{TD Error: } \delta_t}
    \end{equation}

    where $\tilde{c}_t = \arg\max_{c} \hat{Q}_{t-1}(c, s_t)$ is the optimal action in the new state under current beliefs. The Kalman gain $\alpha^E_t(c,s)$ is defined by the GP covariance:
    \begin{equation}
        \alpha^E_t(c,s) = \frac{\text{Cov}(\delta_t, Q^*(c,s) \mid \mathcal{I}_{t-1})}{\text{Var}(\delta_t \mid \mathcal{I}_{t-1})}
    \end{equation}
    The posterior covariance $\hat{\Sigma}^E_t$ is updated to reflect the reduction in uncertainty:
    \begin{equation}
        \hat{\Sigma}^E_t(c,s,c',s') = \hat{\Sigma}_{t-1}(c,s,c',s') - \alpha^E_t(c,s) \text{Cov}(\delta_t, Q^*(c',s') \mid \mathcal{I}_{t-1})
    \end{equation}
\end{definition}

\paragraph{Noise in Experience Signal} Because the experience signal is based on a single realized transition, it is subject
to noise from the stochastic transition function $\mathcal{F}$. Agents may
misattribute actions to outcomes due to random shocks. This can then cause posterior updates
that updates upwards or downards incorrectly.

\subsubsection{An Example of Experience-Based Learning in a Two-Period Model}

Consider a two-period model in which an agent is endowed with initial wealth $s_1$ and 
chooses consumption $c_1 \in [0, s_1]$, saving the remainder at gross interest rate $R$.
Second-period wealth is $s_2 = R(s_1 - c_1) + \varepsilon$, where
$\varepsilon \sim \mathcal{N}(0, \sigma_\varepsilon^2)$ is a mean-zero shock to the 
transition. The agent has quadratic utility $u(c) = c - \frac{1}{2}c^2$ and
discounts at rate $\beta$. In the terminal period the agent 
consumes all wealth, so $c_2 = s_2$. The agent's problem is:

\begin{equation}
    \max_{c_1 \in [0, s_1]} \left[c_1 - \frac{1}{2}c_1^2\right]
    + \beta \, \mathbb{E}\left[s_2 - \frac{1}{2}s_2^2\right]
\end{equation}

We set $\beta R = 1$, under which certainty equivalence holds and the rational agent smooths consumption
equally across periods. The optimal first-period consumption is:
\begin{equation}
    c_1^* = \frac{R \, s_1}{1 + R}
\end{equation}

In our decision framework this solution implies the following state-value and action-value functions:

\begin{align*}
    Q^*(c_1, s_1) &= c_1 - \frac{1}{2}c_1^2 + \beta \left[s_2 - \frac{1}{2}s_2^2 + \sigma_\varepsilon^2\right] \\
    V^*(s_1) &= Q^*(c_1^*, s_1), \quad V^*(s_2) = s_2 - \frac{1}{2}s_2^2
\end{align*}

We next consider an agent with the same utility and budget set who does not know the optimal policy or
value functions. The agent holds a GP prior over
$Q^*$:
\begin{equation}
  Q^*(s,c) \sim GP\left(
    \hat{Q}_0(s,c), \;
    k\left((s,c),(s',c')\right)
  \right)
\end{equation}
with RBF kernel:
\begin{equation}
  k\left((s,c),(s',c')\right) = \sigma_0^2
    \exp\left(
      -\psi_s(s - s')^2 - \psi_c(c - c')^2
    \right)
\end{equation}
We center the prior mean on the true action value function, $\hat{Q}_0 = Q^*$, so that any posterior
deviation from the truth is driven entirely by noise in the experience signal rather than prior
misspecification. We assume the agent chooses optimal action $c_1^*$ in the first period, 
but is then subject to some transmission shock $\varepsilon$. Their temporal difference 
signal is then:

\begin{align*}
    \delta &= u(c_1^*, s_1) + \beta \hat{Q}_0(s_2, s_2) - \hat{Q}_0(c_1^*, s_1) \\
    &= c_1^* - \tfrac{1}{2}(c_1^*)^2 + \beta \hat{Q}_0(c_1^* + \varepsilon, c_1^* + \varepsilon) - \hat{Q}_0(c_1^*, s_1) \\
    &= c_1^* - \tfrac{1}{2}(c_1^*)^2 + \beta \left[c_1^* + \varepsilon - \tfrac{1}{2}(c_1^* + \varepsilon)^2\right] - \left[c_1^* - \tfrac{1}{2}(c_1^*)^2 + \beta \mathbb{E}\left[s_2 - \tfrac{1}{2}s_2^2\right]\right] \\
    &= c_1^* - \tfrac{1}{2}(c_1^*)^2 + \beta \left[c_1^* + \varepsilon - \tfrac{1}{2}(c_1^* + \varepsilon)^2\right] - c_1^* + \tfrac{1}{2}(c_1^*)^2 - \beta \left[c_1^* - \tfrac{1}{2}(c_1^*)^2 + \sigma_\varepsilon^2\right] \\
    &= \beta \left[\varepsilon (1 - c_1^*) - \tfrac{1}{2}\varepsilon^2 + \tfrac{1}{2}\sigma_\varepsilon^2\right]
\end{align*}

The Kalman gain $\alpha^E$ is described by the covariance between the TD error $\delta$ 
and the action value function $Q^*$:

\begin{align*}
  \alpha^E(s,c) &= \frac{
    \text{Cov}(\delta, Q^*(s,c))
  }{
    \text{Var}(\delta)
  } \\[8pt]
  &= \frac{
    \exp\left[
      -\psi_s(s - s_1)^2
      - \psi_c(c - c_1^*)^2
    \right]
    - \beta \exp\left[
      -\psi_s(\varepsilon + c_1^* - s)^2
      - \psi_c(\varepsilon + c_1^* - c)^2
    \right]
  }{
    1 + \beta^2
    - 2\beta \exp\left[
      -\psi_c \varepsilon^2
      - \psi_s \frac{
        (\varepsilon + \varepsilon R - s_1)^2
      }{(1+R)^2}
    \right]
  }
\end{align*}

Meaning for any point $(s,c)$ in the state-action space, the agent's difference 
between the posterior mean $\hat{Q}^E(s,c)$ and the prior mean $\hat{Q}_0(s,c)$ is:

\begin{align*}
  \hat{Q}^E(s,c) - \hat{Q}_0(s,c)
  &= \alpha^E(s,c) \cdot \delta \\[10pt]
  &= \frac{
    \exp\left[
      -\psi_s(s - s_1)^2
      - \psi_c(c - c_1^*)^2
    \right]
    - \beta \exp\left[
      -\psi_s(\varepsilon + c_1^* - s)^2
      - \psi_c(\varepsilon + c_1^* - c)^2
    \right]
  }{
    1 + \beta^2
    - 2\beta \exp\left[
      -\psi_c \varepsilon^2
      - \psi_s \frac{
        (\varepsilon + \varepsilon R - s_1)^2
      }{(1+R)^2}
    \right]
  } \\
  &\qquad \times \;
    \beta \left[\varepsilon (1 - c_1^*) - \tfrac{1}{2}\varepsilon^2 + \tfrac{1}{2}\sigma_\varepsilon^2\right]
\end{align*}

The difference between the posterior and prior mean at the optimal action $c_1^*$ is
contaminated by noise from the transition shock $\varepsilon$. Given a risk-averse agent, 
if $\varepsilon$ is negative, then the agent experiences a worse outcome than expected, leading to a 
negative TD error and a downward update in beliefs. This is is shown in Figure \ref{fig:update_vs_eps},
which plots the posterior distortion at the period-1 state-action pair $(s_1, c_1^*)$ against the shock $\varepsilon$. 

%insert pfg figure
\begin{figure} \label{fig:update_vs_eps}
\centering
\resizebox{\textwidth}{!}{
  \input{../figures/update_vs_eps.pgf}
}
\caption[Posterior Distortion]{
  \textbf{Posterior Distortion at $(s_1, c_1^*)$
  vs.\ Shock. ($\varepsilon$)}
The agent's posterior mean at starting state $s_1$ and optimal consumption $c_1^*$, $\hat{Q}^E(s_1,c_1^*)$, after the experience-based update, compared to the prior mean $\hat{Q}_0(s_1,c_1^*)$.}
\end{figure}

The agent misattributes the shock to the action. Figure \ref{fig:update_surface} shows the Kalman
gain across the state-action space for three shock realizations. Under a $-2\sigma$ shock (left), the
agent revises beliefs sharply downward near $(s_1, c_1^*)$, concluding the optimal action was
worse than expected. Under a $+2\sigma$ shock (right), the same action appears better than
expected. Comparing the two panels, the downward revision under the negative shock is larger in
magnitude than the upward revision under the equally-sized positive shock. Under no shock (center), the agent
still revises slightly upward: the realized outcome at $\varepsilon = 0$ exceeds what $Q^*$
predicted in expectation, because $Q^*$ prices in the variance of future shocks that did not
materialize, producing a small positive TD error of $\beta \tfrac{1}{2}\sigma_\varepsilon^2$.

%insert pfg figure
\begin{figure} \label{fig:update_surface}
\centering
\resizebox{\textwidth}{!}{\input{../figures/update_surface.pgf}}
\caption[Posterior Distortion Surface]{
  \textbf{Posterior Distortion Surface.} Kalman gain $\alpha^E(s,c)$ across the state-action space for three shock realizations.
  Green indicates upward revision, red indicates downward. The light blue marker denotes the
  period-1 state-action pair $(s_1, c_1^*)$; the dark blue marker denotes the period-2
  point $(s_2, s_2)$. Colors are displayed on a $\mathrm{sgn}(\cdot)\sqrt{\ln(1+|\cdot|)}$ scale,
  which linearizes the RBF kernel's exponential-quadratic decay in distance, so that the color
  gradient is approximately proportional to distance from the observation points. Panel
  titles report the raw error $\delta$; the full posterior update at any point is
  $\alpha^E(s,c) \cdot \delta$.}
\end{figure}

\subsubsection{An Example of Experience-Based Learning in Infinite-Horizon Model}

We extend the two-period example to an infinite-horizon setting with adjustments for
tractability. Consider an agent endowed with initial wealth $s_1$ who chooses a consumption
path $\{c_t\}_{t=1}^\infty$. The agent has log utility $u(c) = \ln(c)$ and faces a multiplicative
wealth transition $s_{t+1} = R(s_t - c_t) \cdot e^{\varepsilon_t}$ with $\varepsilon_t \sim
\mathcal{N}(0, \sigma_\varepsilon^2)$ i.i.d. The agent's problem is:

\begin{equation}
  \max_{\{c_t\}_{t=1}^\infty}
    \sum_{t=1}^\infty \beta^{t-1}
    \mathbb{E}\left[\ln(c_t)\right]
\end{equation}

The optimal policy is $c_t^* = (1-\beta)s_t$
for all $t$, independent of the interest rate $R$.
Under this policy, the true action value function
and state value function are log-linear in their
arguments. Defining the constant
$\xi(\beta,R) \equiv \frac{\ln(1-\beta)}{1-\beta}
+ \frac{\beta\ln(R\beta)}{(1-\beta)^2}$:
\begin{align}
  Q^*(s,c) &= \ln(c)
    + \frac{\beta}{1-\beta}\ln(R(s-c))
    + \xi(\beta, R) \\
  V^*(s) &= \frac{1}{1-\beta}\ln s
    + \xi(\beta, R)
\end{align}

\paragraph{TD Error.}
Suppose the agent holds a GP prior centered on the
truth, $\hat{Q}_0 = Q^*$, and chooses the optimal
action $c_1^* = (1-\beta)s_1$ at initial wealth
$s_1$. Upon transitioning to
$s_2 = \beta R s_1 e^{\varepsilon}$, the temporal difference reduces to a simple function of the shock $\varepsilon$.

\begin{equation}
  \delta
    = \frac{\beta}{1-\beta}
      \left[
        \ln(\beta R s_1)
        - \ln(\beta R s_1 e^{\varepsilon})
      \right]
    = -\frac{\beta}{1-\beta}\,\varepsilon
\end{equation}

The temporal difference is linear in the shock with no
higher-order terms. Its variance is $\mathrm{Var}(\delta) = \left(\tfrac{\beta}{1-\beta}\right)^2
\sigma_\varepsilon^2$, scaling with the squared infinite-horizon discount multiplier. Because
the optimal policy and value functions are stationary, this result holds at any period: an
agent at state $s_{t-1}$ who takes optimal action $c_{t-1}^* = (1-\beta)s_{t-1}$ and transitions to
$s_t = \beta R s_{t-1} e^{\varepsilon_t}$ experiences the same temporal difference
$\delta_t = -\frac{\beta}{1-\beta}\varepsilon_t$,
independent of wealth.

\paragraph{Posterior Update.} To respect the multiplicative nature of the wealth transition, the agent's similarity 
measure must be scale-invariant. The natural choice is the \textit{Log-Laplace kernel}.
\begin{equation}
k\big((s,c), (s',c')\big) = \sigma_0^2 \exp\left( -\psi_s \left|\ln\frac{s}{s'}\right| - \psi_c \left|\ln\frac{c}{c'}\right| \right)
\end{equation}
The form captures the correct geometry, but the absolute values require piecewise 
integration. However, we can derive a global closed-form solution by 
exploiting the identity $\exp(-|z|)=\min(e^z, e^{-z})$. Applied to log-ratios, this 
transforms the kernel into the product of minimum ratios.
\begin{equation}
\exp\left(-\psi \left|\ln\frac{x}{y}\right|\right) = \left[ \min\left(\frac{x}{y}, \frac{y}{x}\right) \right]^\psi
\end{equation}
This allows us to evaluate the kernel covariance analytically across the entire state space.
Let $\rho(\varepsilon)$ denote the realized growth factor of wealth from $t-1$ to period $t$.
Substituting the optimal policy into the transition equation $s_t = R(s_{t-1} - c_{t-1}^*) e^\varepsilon$ yields:
\begin{equation}
\rho(\varepsilon) \equiv \frac{s_t}{s_{t-1}} = \frac{c_t^*}{c_{t-1}^*} = \beta R e^\varepsilon
\end{equation}
Reiterating that the agent's prior is centered on the truth, the decision point is 
$(s_{t-1}, c_{t-1}^*)$ and the outcome point is 
$(s_t, c_t^*) = (s_{t-1}\rho(\varepsilon), c_{t-1}^*\rho(\varepsilon))$. 
Using the simplified min-ratio form, the update is:
\begin{equation}
\hat{Q}^E(s,c) - \hat{Q}_0(s,c) =
\left( -\frac{\beta \varepsilon}{1-\beta} \right)
\cdot
\frac{ \mathcal{K}_{t-1}(s,c) - \beta \mathcal{K}_t(s,c, \varepsilon) }{ 1 + \beta^2 - 2\beta \left[ \min\left(\rho(\varepsilon), \rho(\varepsilon)^{-1}\right) \right]^{\psi_s+\psi_c} }
\end{equation}
where the spatial decay terms $\mathcal{K}_{t-1}$ (anchored at the decision) and $\mathcal{K}_t $ (anchored at the outcome) are given by:
\begin{align}
\mathcal{K}_{t-1}(s,c) &= \left[ \min\left(\frac{s}{s_{t-1}}, \frac{s_{t-1}}{s}\right) \right]^{\psi_s} \left[ \min\left(\frac{c}{c_{t-1}^*}, \frac{c_{t-1}^*}{c}\right) \right]^{\psi_c} \\
\mathcal{K}_t(s,c, \varepsilon) &= \left[ \min\left(\frac{s}{s_t}, \frac{s_t}{s}\right) \right]^{\psi_s} \left[ \min\left(\frac{c}{c_t^*}, \frac{c_t^*}{c}\right) \right]^{\psi_c}
\end{align}

% ============================================================
% Kernel Evaluations at Decision and Outcome Points
% ============================================================

\paragraph{Kernel Evaluations.}
Let $\psi \equiv \psi_s + \psi_c$ and define the \textit{min-ratio kernel overlap}:
\begin{equation}
m(\varepsilon) \equiv \left[\min\!\left(\rho(\varepsilon),\; \rho(\varepsilon)^{-1}\right)\right]^{\psi}
\end{equation}
where $\rho(\varepsilon) = \beta R e^\varepsilon$. Both savings and consumption scale by the common factor $\rho(\varepsilon)$. Each kernel evaluation reduces to a power of $\rho$ alone.
\begin{align}
\mathcal{K}_{t-1}(s_{t-1}, c_{t-1}^*) &= 1, &
\mathcal{K}_t(s_{t-1}, c_{t-1}^*, \varepsilon) &= m(\varepsilon), \notag\\
\mathcal{K}_{t-1}(s_t, c_t^*) &= m(\varepsilon), &
\mathcal{K}_t(s_t, c_t^*, \varepsilon) &= 1, \notag\\[6pt]
%
\hat{Q}^E(s_{t-1}, c_{t-1}^*) - \hat{Q}_0(s_{t-1}, c_{t-1}^*)
&= \frac{-\beta\varepsilon}{(1-\beta)}
\cdot \frac{1 - \beta\, m(\varepsilon)}{1 + \beta^2 - 2\beta\, m(\varepsilon)}, \\[6pt]
%
\hat{Q}^E(s_t, c_t^*) - \hat{Q}_0(s_t, c_t^*)
&= \frac{-\beta\varepsilon}{(1-\beta)}
\cdot \frac{m(\varepsilon) - \beta}{1 + \beta^2 - 2\beta\, m(\varepsilon)}.
\end{align}
The two updates share a common denominator and prefactor. 
\begin{equation}
m(\varepsilon) =
\begin{cases}
(\beta R e^\varepsilon)^{\psi} & \text{if } \varepsilon < \varepsilon^* \quad (\rho < 1), \\[4pt]
(\beta R e^\varepsilon)^{-\psi} & \text{if } \varepsilon \geq \varepsilon^* \quad (\rho \geq 1).
\end{cases}
\end{equation}




% We illustrate the experience-based learning mechanism with a tractable consumption-savings example.

% \paragraph{Environment.} Consider an agent with quadratic utility $u(c) = c - \frac{\alpha}{2}c^2$ facing a linear wealth transition $s_{t+1} = R(s_t - c_t) + y_t$, where $y_t \in \{y_L, y_M, y_H\}$ is income drawn uniformly, so $\mathbb{E}[y_t] = y_M$. We impose $\beta R = 1$, which implies the optimal policy is to smooth consumption: $c^* = \frac{R-1}{R}s + y_M$.

% \paragraph{Functional Form of $Q^*$.} With quadratic utility and linear transitions, the Bellman equation admits a polynomial solution. Therefore:
% \begin{equation}
%     Q^*(s,c) = \theta_0 + \theta_1 s + \theta_2 c + \theta_3 c^2 + \theta_4 sc + \theta_5 s^2
% \end{equation}
% where the coefficients $\boldsymbol{\theta} = (\theta_0, \theta_1, \theta_2, \theta_3, \theta_4, \theta_5)^\top$ depend on the primitives $\alpha, \beta, R, y_M, \text{Var}(y)$.

% \paragraph{Agent's Prior.} The agent knows $Q^*$ is polynomial in $(s,c)$ but does not know the coefficients. They hold a Gaussian prior:
% \begin{equation}
%     \boldsymbol{\theta} \sim \mathcal{N}(\boldsymbol{\mu}_0, \boldsymbol{\Sigma}_0)
% \end{equation}
% Define the feature vector $\boldsymbol{\phi}(s,c) = (1, s, c, c^2, sc, s^2)^\top$, so that $Q^*(s,c) = \boldsymbol{\phi}(s,c)^\top \boldsymbol{\theta}$.

% \paragraph{Greedy Policy.} Given current beliefs $\hat{\boldsymbol{\mu}}_{t-1}$, the agent chooses:
% \begin{equation}
%     c_t = \arg\max_c \, \boldsymbol{\phi}(s_t, c)^\top \hat{\boldsymbol{\mu}}_{t-1}
% \end{equation}

% \paragraph{Bayesian Updating.} The agent observes utility $\eta^E_t = u(c_{t-1})$ and interprets it via the TD approximation:
% \begin{equation}
%     \eta^E_t = Q^*(s_{t-1}, c_{t-1}) - \beta Q^*(s_t, \tilde{c}_t) = \boldsymbol{x}_t^\top \boldsymbol{\theta}
% \end{equation}
% where $\boldsymbol{x}_t = \boldsymbol{\phi}(s_{t-1}, c_{t-1}) - \beta \boldsymbol{\phi}(s_t, \tilde{c}_t)$ and $\tilde{c}_t$ is the greedy action at $s_t$. The posterior update is:
% \begin{align}
%     \hat{\boldsymbol{\mu}}_t &= \hat{\boldsymbol{\mu}}_{t-1} + \boldsymbol{K}_t \left(\eta^E_t - \boldsymbol{x}_t^\top \hat{\boldsymbol{\mu}}_{t-1}\right) \\
%     \hat{\boldsymbol{\Sigma}}_t &= \hat{\boldsymbol{\Sigma}}_{t-1} - \boldsymbol{K}_t \boldsymbol{x}_t^\top \hat{\boldsymbol{\Sigma}}_{t-1}
% \end{align}
% where the Kalman gain is:
% \begin{equation}
%     \boldsymbol{K}_t = \frac{\hat{\boldsymbol{\Sigma}}_{t-1} \boldsymbol{x}_t}{\boldsymbol{x}_t^\top \hat{\boldsymbol{\Sigma}}_{t-1} \boldsymbol{x}_t}
% \end{equation}

% \paragraph{Misattribution from Income Shocks.} Suppose the agent experiences income sequence $y_L, y_M, y_M, y_H, y_M, y_M$. In period 1, income $y_L < y_M$ realizes. The agent arrives at state $s_1 = R(s_0 - c_0) + y_L$, which is lower than expected. The TD error $\eta^E_1 - \boldsymbol{x}_1^\top \hat{\boldsymbol{\mu}}_0$ is negative: the agent received less than predicted. They update $\hat{\boldsymbol{\mu}}_1$ downward, concluding that consuming $c_0$ at wealth $s_0$ was worse than previously believed. But the bad outcome came from unlucky income, not from the consumption choice. The agent misattributes luck to policy.

% When $y_H$ realizes in period 4, the opposite occurs: the agent arrives at a higher state than expected, the TD error is positive, and they update upward---again misattributing the income shock to their action. Over the sequence, beliefs oscillate with income realizations rather than converging smoothly to the truth.

\subsubsection{Learning Through Reasoning}

We suppose that the agent can produce mental simulations and evaluate counterfactual actions, developing a 
noisy vector of signals $\boldsymbol{\eta}^R_t$ that provide information about the 
action value function. To formalize this, fix the current state $s_t$ and consider 
the vector of action values across all possible consumption levels:
\begin{equation}
    \boldsymbol{Q}^*(s_t) = (Q^*(s_t, c_1), Q^*(s_t, c_2), \ldots, Q^*(s_t, c_n))^\top
\end{equation}
By the GP assumption, this vector is jointly Gaussian with mean 
$\hat{\boldsymbol{Q}}_{t-1}(s_t)$ and covariance $\hat{\boldsymbol{\Sigma}}_{t-1}(s_t)$ 
derived from the kernel. We then suppose the agent can generate noisy observations of
linear combinations of this vector, which represent mental simulations of different actions.
The noise covariance captures the precision of the agent's mental model. The higher the precision,
the more accurate the simulations, but the greater the cognitive cost. In the limit,
the agent is able to perfectly evaluate all actions, equivalent to knowing $Q^*$ exactly.

\begin{definition}[Reasoning-Based Learning Update]
    Following the experience-based update, the agent performs internal mental 
    simulations to refine their beliefs. At the current state $s_t$, the agent 
    generates a vector of noisy reasoning signals $\boldsymbol{\eta}^R_t$, 
    representing hypothesized action-values. These signals are modeled as:
    \begin{equation}
        \boldsymbol{\eta}^R_t = \boldsymbol{\Omega}'_t \boldsymbol{Q}^*(s_t) + \boldsymbol{\epsilon}_{\eta, t}, \quad \boldsymbol{\epsilon}_{\eta, t} \sim \mathcal{N}(\boldsymbol{0}, \boldsymbol{\Sigma}_{\eta, t})
    \end{equation}
    where $\boldsymbol{\Omega}'_t$ is a $k \times n$ loadings matrix chosen by the 
    agent to determine which action comparisons are evaluated during simulation. 
    Treating these as signals of the true $Q^*$ function, the agent forms reasoned 
    beliefs $(\hat{Q}^R_t, \hat{\Sigma}^R_t)$ via a Bayesian update on the intermediate 
    beliefs $(\hat{Q}^E_t, \hat{\Sigma}^E_t)$:
    \begin{align}
        \hat{Q}^R_t(c,s) &= \hat{Q}^E_t(c,s) + \boldsymbol{\alpha}^R_t(c,s)^\top \left[\boldsymbol{\eta}^R_t - \boldsymbol{\Omega}'_t \hat{\boldsymbol{Q}}^E_t(s_t)\right] \\
        \boldsymbol{\alpha}^R_t(c,s) &= \text{Cov}(\boldsymbol{\eta}^R_t, Q^*(c,s) \mid \mathcal{I}^E_t) \left[\text{Var}(\boldsymbol{\eta}^R_t \mid \mathcal{I}^E_t)\right]^{-1}
    \end{align}
    The posterior covariance $\hat{\Sigma}^R_t$ is updated accordingly to reflect the reduction in uncertainty from mental simulation:
    \begin{equation}
        \hat{\Sigma}^R_t(c,s,c',s') = \hat{\Sigma}^E_t(c,s,c',s') - \boldsymbol{\alpha}^R_t(c,s)^\top \text{Cov}(\boldsymbol{\eta}^R_t, Q^*(c',s') \mid \mathcal{I}^E_t)
    \end{equation}
\end{definition}

\subsection{Action Selection}

At each period $t$, the agent goes through the set of steps in Definition \ref{def:learning_cycle}. They first perform a bayesian
update of their beliefs based on realized utility from the previous period: \textit{experience-based learning}.
They then must choose how much reasoning to do to optimally update their beliefs: \textit{reasoning decision}.
At each step, the agent faces two tasks: one to determine optimal policy given current beliefs, and another to determine
the optimal reasoning strategy to improve those beliefs.

\subsubsection{The Optimal Policy Decision}

The agent's action decision is governed by their policy function $\pi_t(c | s_t)$. This policy is derived from maximizing expected payoffs subject to exploration constraints. The agent solves the following problem based on their current beliefs $\hat{Q}_t$:

\begin{align}
    \max_{\pi_t(\cdot | s_t)} \quad & \sum_c \hat{Q}_t(c, s_t) \pi_t(c | s_t) \\
    \text{s.t.} \quad & \sum_c \pi_t(c | s_t) = 1 \\
    & -\sum_c \pi_t(c | s_t) \ln \pi_t(c | s_t) \geq h \sum_c \sigma_t(c, s_t) \label{eq:entropy_constraint}
\end{align}

The exploration constraint is motivated by the principle of 
\textit{uncertainty-driven exploration}. In standard maximum entropy reinforcement 
learning, the entropy floor is often a fixed parameter. Here, we specify the lower bound 
as proportional to the agent's total posterior uncertainty, $\sum \sigma_t(c, s_t)$. 
Experimentation is valuable only to the extent that the agent is uncertain about the true payoffs. 
When uncertainty is high, the constraint tightens, forcing the agent to spread 
probability mass across actions to learn; as uncertainty vanishes, the constraint 
relaxes, allowing the agent to converge toward the greedy policy.

\begin{proposition}[Optimal Action Policy] \label{prop:optimal_policy}
    The solution to the constrained maximization problem yields a state-dependent softmax distribution:
    \begin{equation}
        \pi_t(c | s_t) = \begin{cases} 
        \dfrac{\exp\left(\hat{Q}_t(c, s_t) / \delta_t\right)}{\sum_{c'} \exp\left(\hat{Q}_t(c', s_t) / \delta_t\right)} & \text{if } \delta_t > 0 \\[3ex] 
        \mathbb{I}\left(c = \arg\max_{c'} \hat{Q}_t(c', s_t)\right) & \text{if } \delta_t = 0 
        \end{cases}
    \end{equation}
    where $\delta_t \geq 0$ is the Lagrange multiplier associated with the entropy constraint.
    The multiplier $\delta_t$ acts as an endogenous temperature parameter. When 
uncertainty is high, $\delta_t$ is large, flattening the policy distribution.
\end{proposition}

 Note
that this optimization is performed twice within the period $t$ learning cycle. First, 
it is solved using the intermediate experience-based beliefs $\hat{Q}^E_t$ to form draft policy $\pi^E_t$,
which is used in the reasoning process to evaluate expected payoffs from mental simulation.
Second, after reasoning signals are realized, it is solved using the final reasoned beliefs $\hat{Q}^R_t$ 
to determine the $\pi^R_t$ for action execution.

\subsubsection{The Reasoning Decision}\label{sec:reasoning-decision}

Reducing uncertainty allows the agent to tighten their priors over the action value function,
and thus improve their policy. Because they are endowed with some reasoning capacity, they
will naturally use it to improve their beliefs. However, reasoning is costly.
The agents face a rational inattention problem in which they must optimally allocate their limited
cognitive resources to improve their beliefs. They choose how many simulations to run
and the precision of those simulations, balancing the expected improvement in payoffs
against the cognitive cost of reasoning.

The agent chooses the simulations to run ($\boldsymbol{\Omega}_t$) and the precision of 
those simulations ($\boldsymbol{\Sigma}_{\eta, t}$). The reasoning signal takes the 
form $\eta^R_t = \boldsymbol{\Omega}_t Q^* + \varepsilon$ with $\varepsilon 
\sim \mathcal{N}(0, \boldsymbol{\Sigma}_{\eta,t})$, where rows of $\boldsymbol{\Omega}_t$ 
specify which linear combinations of $Q^*(c,s)$ to reason about, and $\boldsymbol{\Sigma}_{\eta,t}$ 
governs the noise in those comntemplations. The agent maximizes expected payoff net of reasoning costs:

\begin{equation}
\max_{\boldsymbol{\Omega}_t, \boldsymbol{\Sigma}_{\eta, t}} \mathbb{E}\left[ \sum_c \hat{Q}^R_t(c, s_t) \cdot \pi_t(c | s_t) \mid \mathcal{I}^E_t \right] - \kappa \ln\left[\frac{|\hat{\boldsymbol{\Sigma}}^E_t|}{|\hat{\boldsymbol{\Sigma}}^R_t|}\right]
\end{equation}
The cost term is the information gain from reasoning, measured as the entropy reduction in beliefs about $Q^*$. The lower the posterior covariance $\hat{\boldsymbol{\Sigma}}^R_t$, the more precise the beliefs, and the higher the cost; in the limit the cost of perfect information becomes infinite. The parameter $\kappa$ governs the marginal cost of information.

\begin{proposition}[Optimal Reasoning Allocation]
    Let $\hat{\boldsymbol{\Sigma}}^E_t = V \Lambda^E V'$ be the eigendecomposition of the prior covariance, where $\Lambda^E = \mathrm{diag}(\lambda_1^E, \ldots, \lambda_n^E)$ with $\lambda_1^E \geq \cdots \geq \lambda_n^E > 0$. The optimal posterior covariance is:
    \begin{equation}
        \hat{\boldsymbol{\Sigma}}^R_t = V \Lambda^* V'
    \end{equation}
    where $\Lambda^* = \mathrm{diag}(\lambda_1^*, \ldots, \lambda_n^*)$ with:
    \begin{equation}
        \lambda_i^* = \min\left\{ \lambda_i^E, \; \frac{\kappa}{h \delta^E_t} \right\}
    \end{equation}
    and $\delta^E_t$ is the Lagrange multiplier from the draft policy problem solved under prior beliefs $\hat{Q}^E_t$. The optimal signal structure sets $\boldsymbol{\Omega}^*_t = V'$ and $\boldsymbol{\Sigma}^*_{\eta,t} = \mathrm{diag}(\sigma_1^2, \ldots, \sigma_n^2)$ with:
    \begin{equation}
        \sigma_i^2 = \begin{cases} 
            \dfrac{\lambda_i^E \lambda_i^*}{\lambda_i^E - \lambda_i^*} & \text{if } \lambda_i^E > \dfrac{\kappa}{h\delta^E_t} \\[2ex]
            \infty & \text{otherwise}
        \end{cases}
    \end{equation}
\end{proposition}

\paragraph{Reverse Water Filling.} The threshold $\kappa/(h\delta^E_t)$ acts as a ``water level.'' 
Eigendirections with prior variance exceeding this level are reasoned about until 
variance equals the threshold. The agent allocates cognitive effort to their largest 
uncertainties first, stopping when the marginal cost of further precision ($\kappa$) 
exceeds its marginal value ($h\delta^E_t$).

The reasoning decision must be made before the reasoning signal $\eta^R_t$ is realized, 
so the agent cannot condition on the final posterior $\hat{Q}^R_t$. Instead, they 
evaluate the marginal value of precision using $\delta^E_t$—the shadow price of the 
entropy constraint from the draft policy problem solved under experience-based beliefs 
$\hat{Q}^E_t$. This shadow price measures how much the agent would gain from being 
allowed to exploit more (randomize less). Under certainty equivalence, the agent 
commits to a reasoning intensity as if the posterior mean will equal the prior mean, 
making $\delta^E_t$ a sufficient statistic for how valuable reducing uncertainty would be.